\textbf{Phần 1. Trắc nghiệm lựa chọn}
\setcounterpageref{ex}{0}

\begin{ex} %[5D12N72]
$1\text{ m}^3$ bằng bao nhiêu đề-xi-mét khối?
\choice
{$10\text{ dm}^3$}
{$100\text{ dm}^3$}
{\True $1,000\text{ dm}^3$}
{$10,000\text{ dm}^3$}
\loigiai{
$1\text{ m}^3 = 1,000\text{ dm}^3$.\\
\textbf{Đáp án của câu hỏi là C}
}
\end{ex}

\begin{ex} %[5D12N72]
$3\text{ tấn } 50\text{ kg}$ viết dưới dạng số thập phân có đơn vị là tấn là
\choice
{$3{,}5\text{ tấn}$}
{$3{,}005\text{ tấn}$}
{$30{,}5\text{ tấn}$}
{\True $3{,}05\text{ tấn}$}
\loigiai{
$3\text{ tấn } 50\text{ kg} = 3 + \dfrac{50}{1000} = 3{,}05\text{ tấn}$.\\
\textbf{Đáp án của câu hỏi là D}
}
\end{ex}

\begin{ex} %[5D12H72]
$2{,}4\text{ ha}$ bằng bao nhiêu mét vuông?
\choice
{$240\text{ m}^2$}
{\True $24,000\text{ m}^2$}
{$2,400\text{ m}^2$}
{$240,000\text{ m}^2$}
\loigiai{
$1\text{ ha} = 10,000\text{ m}^2 \Rightarrow 2{,}4\text{ ha} = 24,000\text{ m}^2$.\\
\textbf{Đáp án của câu hỏi là B}
}
\end{ex}

\begin{ex} %[5D12H72]
$1$ giờ $45$ phút bằng bao nhiêu giờ?
\choice
{\True $1{,}75\text{ giờ}$}
{$1{,}45\text{ giờ}$}
{$1{,}25\text{ giờ}$}
{$1{,}5\text{ giờ}$}
\loigiai{
$1 + 45 : 60 = 1{,}75\text{ giờ}$.\\
\textbf{Đáp án của câu hỏi là A}
}
\end{ex}

\begin{ex} %[5D12H72]
$450\text{ dm}^3$ bằng bao nhiêu mét khối?
\choice
{$4{,}5\text{ m}^3$}
{$0{,}045\text{ m}^3$}
{\True $0{,}45\text{ m}^3$}
{$45\text{ m}^3$}
\loigiai{
$450 : 1,000 = 0{,}45\text{ m}^3$.\\
\textbf{Đáp án của câu hỏi là C}
}
\end{ex}

\begin{ex} %[5D12H72]
Điền số thích hợp vào chỗ chấm: $5\text{ m } 8\text{ cm} = \dots\text{ m}$.
\choice
{\True $5{,}08$}
{$5{,}8$}
{$5{,}008$}
{$58$}
\loigiai{
$5 + 8 : 100 = 5{,}08\text{ m}$.\\
\textbf{Đáp án của câu hỏi là A}
}
\end{ex}

\begin{ex} %[5D12V72]
Một thửa ruộng thu hoạch được $2{,}4\text{ tấn}$ thóc. Người ta đóng đều vào $48$ bao. Mỗi bao nặng bao nhiêu ki-lô-gam thóc?
\choice
{$40\text{ kg}$}
{$60\text{ kg}$}
{$55\text{ kg}$}
{\True $50\text{ kg}$}
\loigiai{
$2{,}4\text{ tấn} = 2,400\text{ kg}$.\\
Mỗi bao: $2,400 : 48 = 50\text{ (kg)}$.\\
\textbf{Đáp án của câu hỏi là D}
}
\end{ex}

\begin{ex} %[5D12V72]
Một mảnh đất có diện tích $0{,}15\text{ ha}$. Người ta dành $40\%$ diện tích để xây nhà. Diện tích đất xây nhà là
\choice
{$500\text{ m}^2$}
{\True $600\text{ m}^2$}
{$750\text{ m}^2$}
{$60\text{ m}^2$}
\loigiai{
$0{,}15\text{ ha} = 1,500\text{ m}^2$.\\
Diện tích xây nhà: $1,500 \times 40 : 100 = 600\text{ (m}^2\text{)}$.\\
\textbf{Đáp án của câu hỏi là B}
}
\end{ex}

\begin{ex} %[5D12V72]
Một bể nước chứa được $1{,}8\text{ m}^3$ nước. Hiện tại trong bể đã có $600\text{ l}$ nước. Cần đổ thêm bao nhiêu lít nước nữa để đầy bể?
\choice
{\True $1,200\text{ l}$}
{$1,000\text{ l}$}
{$1,400\text{ l}$}
{$800\text{ l}$}
\loigiai{
$1{,}8\text{ m}^3 = 1,800\text{ l}$.\\
Cần đổ thêm: $1,800 - 600 = 1,200\text{ (l)}$.\\
\textbf{Đáp án của câu hỏi là A}
}
\end{ex}

\begin{ex} %[5D12V72]
Sắp xếp các số đo sau theo thứ tự từ bé đến lớn: $3\text{ m}^2 5\text{ dm}^2; 3{,}5\text{ m}^2; 30,050\text{ cm}^2; 3{,}04\text{ m}^2$.
\choice
{$30,050\text{ cm}^2 < 3{,}04\text{ m}^2 < 3\text{ m}^2 5\text{ dm}^2 < 3{,}5\text{ m}^2$}
{$3{,}04\text{ m}^2 < 30,050\text{ cm}^2 < 3\text{ m}^2 5\text{ dm}^2 < 3{,}5\text{ m}^2$}
{\True $3{,}005\text{ m}^2 (30,050\text{ cm}^2) < 3{,}04\text{ m}^2 < 3{,}05\text{ m}^2 (3\text{ m}^2 5\text{ dm}^2) < 3{,}5\text{ m}^2$}
{$3{,}5\text{ m}^2 < 3\text{ m}^2 5\text{ dm}^2 < 3{,}04\text{ m}^2 < 30,050\text{ cm}^2$}
\loigiai{
Đổi về $\text{m}^2$: $30,050\text{ cm}^2 = 3{,}005\text{ m}^2 < 3{,}04\text{ m}^2 < 3{,}05\text{ m}^2 < 3{,}5\text{ m}^2$.\\
\textbf{Đáp án của câu hỏi là C}
}
\end{ex}

\begin{ex} %[5D12V72]
Một cuộn dây thép dài $45\text{ m } 6\text{ dm}$, mỗi mét dây nặng $250\text{ g}$. Cả cuộn dây đó nặng bao nhiêu ki-lô-gam?
\choice
{$11\text{ kg}$}
{\True $11{,}4\text{ kg}$}
{$114\text{ kg}$}
{$1{,}14\text{ kg}$}
\loigiai{
$45\text{ m } 6\text{ dm} = 45{,}6\text{ m}$.\\
Khối lượng: $45{,}6 \times 250 = 11,400\text{ g} = 11{,}4\text{ kg}$.\\
\textbf{Đáp án của câu hỏi là B}
}
\end{ex}

\begin{ex} %[5D12C72]
Một vòi nước chảy vào bể không có nước, mỗi phút chảy được $25\text{ l}$ nước. Sau $1$ giờ $12$ phút thì lượng nước trong bể chiếm $60\%$ dung tích của bể. Dung tích toàn bộ của bể nước đó là
\choice
{$2,500\text{ l}$}
{$3,600\text{ l}$}
{$4,000\text{ l}$}
{\True $3,000\text{ l}$}
\loigiai{
$1$ giờ $12$ phút $= 72$ phút.\\
Lượng nước đã chảy: $72 \times 25 = 1,800\text{ (l)}$.\\
Dung tích bể: $1,800 : 60 \times 100 = 3,000\text{ (l)}$.\\
\textbf{Đáp án của câu hỏi là D}
}
\end{ex}

\textbf{Phần 2. Trắc nghiệm đúng sai}
\setcounterpageref{ex}{0}

\begin{ex} %[5D12H72]
Xét tính đúng sai của các biểu thức đổi đơn vị đo:
\choiceTF
{\True $4\text{ tấn } 8\text{ kg} = 4{,}008\text{ tấn}$.}
{$5\text{ m}^2 6\text{ dm}^2 = 5{,}6\text{ m}^2$.}
{\True $5\text{ m}^2 6\text{ dm}^2 = 5{,}06\text{ m}^2$.}
{\True $2\text{ dm}^3 45\text{ cm}^3 = 2{,}045\text{ dm}^3$.}
\loigiai{
\begin{itemize}
    \item $4\text{ tấn } 8\text{ kg} = 4 + 8 : 1000 = 4{,}008\text{ tấn}$. Đúng.
    \item $5\text{ m}^2 6\text{ dm}^2 = 5 + 6 : 100 = 5{,}06\text{ m}^2 \neq 5{,}6\text{ m}^2$. Ý 2 Sai, Ý 3 Đúng.
    \item $2\text{ dm}^3 45\text{ cm}^3 = 2 + 45 : 1000 = 2{,}045\text{ dm}^3$. Đúng.
\end{itemize}
\begin{itemize}
    \item a) \textbf{Đúng}.
    \item b) \textbf{Sai}.
    \item c) \textbf{Đúng}.
    \item d) \textbf{Đúng}.
\end{itemize}
}
\end{ex}

\begin{ex} %[5D12V72]
Một xe tải chở $3$ chuyến hàng: Chuyến 1 chở $2{,}5\text{ tấn}$, chuyến 2 chở $28\text{ tạ}$, chuyến 3 chở $2,700\text{ kg}$. Xét tính đúng sai:
\choiceTF
{\True Chuyến thứ hai chở nhiều hàng nhất ($2{,}8\text{ tấn}$).}
{\True Chuyến thứ nhất chở ít hàng nhất ($2{,}5\text{ tấn}$).}
{Tổng khối lượng hàng chở trong cả 3 chuyến là $7\text{ tấn}$.}
{\True Tổng khối lượng hàng chở trong cả 3 chuyến là $8\text{ tấn}$.}
\loigiai{
\begin{itemize}
    \item Đổi ra tấn: $2{,}5\text{ tấn}$; $28\text{ tạ} = 2{,}8\text{ tấn}$; $2,700\text{ kg} = 2{,}7\text{ tấn}$.
    \item Chuyến 2 nhiều nhất ($2{,}8\text{ tấn}$). Đúng.
    \item Chuyến 1 ít nhất ($2{,}5\text{ tấn}$). Đúng.
    \item Tổng: $2{,}5 + 2{,}8 + 2{,}7 = 8{,}0\text{ tấn} \neq 7\text{ tấn}$. Ý 3 Sai, Ý 4 Đúng.
\end{itemize}
\begin{itemize}
    \item a) \textbf{Đúng}.
    \item b) \textbf{Đúng}.
    \item c) \textbf{Sai}.
    \item d) \textbf{Đúng}.
\end{itemize}
}
\end{ex}

\textbf{Phần 3. Trả lời ngắn}
\setcounterpageref{ex}{0}

\begin{ex} %[5D12H72]
Đổi $3{,}75\text{ m}^3$ ra đề-xi-mét khối.
\shortans{3750}
\loigiai{
\textbf{Đáp số:} $3750$\\
$3{,}75 \times 1,000 = 3,750\text{ dm}^3$.
}
\end{ex}

\begin{ex} %[5D12H72]
Đổi $5\text{ ha } 400\text{ m}^2$ ra số thập phân có đơn vị là héc-ta.
\shortans{5,04}
\loigiai{
\textbf{Đáp số:} $5{,}04$\\
$5 + 400 : 10,000 = 5{,}04\text{ ha}$.
}
\end{ex}

\begin{ex} %[5D12V72]
Đổi $2$ giờ $36$ phút ra giờ (dùng dấu phẩy nếu có).
\shortans{2,6}
\loigiai{
\textbf{Đáp số:} $2{,}6$\\
$2 + 36 : 60 = 2{,}6\text{ giờ}$.
}
\end{ex}

\begin{ex} %[5D12V72]
Một thửa ruộng hình chữ nhật có diện tích $3,200\text{ m}^2$. Hỏi thửa ruộng đó rộng bao nhiêu héc-ta? (Dùng dấu phẩy nếu có).
\shortans{0,32}
\loigiai{
\textbf{Đáp số:} $0{,}32$\\
$3,200 : 10,000 = 0{,}32\text{ ha}$.
}
\end{ex}

\begin{ex} %[5D12V72]
Một bình nước chứa $2{,}5\text{ l}$ nước. Người ta rót đều vào $10$ chiếc cốc nhỏ. Hỏi mỗi cốc chứa bao nhiêu mi-li-lít nước?
\shortans{250}
\loigiai{
\textbf{Đáp số:} $250$\\
$2,500 : 10 = 250\text{ ml}$.
}
\end{ex}

\begin{ex} %[5D12C72]
Người ta lát một căn phòng hình chữ nhật dài $8\text{ m}$, rộng $6\text{ m}$ bằng những viên gạch hình vuông cạnh $40\text{ cm}$. Giá tiền mỗi viên gạch là $25,000$ đồng. Tính tổng số tiền mua gạch theo nghìn đồng.
\shortans{7500}
\loigiai{
\textbf{Đáp số:} $7500$\\
Diện tích phòng: $8 \times 6 = 48\text{ m}^2 = 480,000\text{ cm}^2$.\\
Diện tích 1 viên gạch: $40 \times 40 = 1,600\text{ cm}^2$.\\
Số viên gạch: $480,000 : 1,600 = 300$ viên.\\
Tổng tiền: $300 \times 25 = 7,500$ nghìn đồng.
}
\end{ex}
